\documentclass[11pt]{article}
\usepackage[T1]{fontenc}
\usepackage[a4paper,margin=1in]{geometry}
\usepackage{amsmath,amssymb,amsthm}
\usepackage[hidelinks]{hyperref}
\usepackage{microtype}

\newtheorem*{question}{Source random-matrix conjecture}
\newtheorem*{answer}{Later theorem}

\setlength{\parindent}{0pt}
\setlength{\parskip}{0.55em}

\title{Literature Status: Tensor-GUE Strong Convergence\\and the Peterson--Thom Conjecture}
\author{Run \texttt{fa\_banach\_001} --- lane 14}
\date{9 August 2026}

\begin{document}
\maketitle

\begin{center}
\fbox{\parbox{0.9\textwidth}{
\textbf{Status: literature already answered.}
Hayes reduces the Peterson--Thom conjecture to strong convergence of two
same-size independent GUE families acting on the two tensor factors.
Belinschi and Capitaine prove exactly this strong convergence in
arXiv:2205.07695, and explicitly note the resulting solution of
Peterson--Thom.
}}
\end{center}

\section{Original conjecture and reduction}

B. Hayes, \emph{A random matrix approach to the Peterson--Thom conjecture},
arXiv:2008.12287~\cite{source}, defines independent $k\times k$ normalized
GUE families $X^{(k)}=(X_i^{(k)})_{i=1}^r$ and
$Y^{(k)}=(Y_i^{(k)})_{i=1}^r$.  Theorem~1.1(iii), on arXiv PDF page~3,
isolates the following statement; Section~4 restates it as a conjecture.

\begin{question}
For every noncommutative polynomial $P$,
\[
 \begin{split}
 &\|P(X^{(k)}\otimes I_k,I_k\otimes Y^{(k)})\|\\
 &\qquad\longrightarrow
 \|P(s\otimes 1,1\otimes s)\|_{C^*(s)\otimes_{\min}C^*(s)}
 \end{split}
\]
in probability, where $s$ is a free semicircular family.
\end{question}

Theorem~1.1 proves that this statement implies positivity of the relevant
1-bounded entropy for every nonamenable subalgebra, which in turn implies the
Peterson--Thom conjecture: every diffuse amenable subalgebra of a free group
factor lies in a unique maximal amenable subalgebra.

\section{Separate later answer}

S. T. Belinschi and M. Capitaine,
\emph{Strong convergence of tensor products of independent G.U.E.
matrices}, arXiv:2205.07695~\cite{answerpaper}, study independent normalized
GUE families $X_N$ and $Y_N$ of the same size.  Their Theorem~1 on arXiv PDF
page~3 states:

\begin{answer}
Almost surely, for every noncommutative polynomial $P$,
\[
 \begin{split}
 &\|P(X_N\otimes I_N,I_N\otimes Y_N)\|\\
 &\qquad\longrightarrow
 \|P(s\otimes 1,1\otimes t)\|_{\min},
 \end{split}
\]
where $s$ and $t$ are free semicircular systems in the two tensor factors.
\end{answer}

This is the source conjecture, with the harmless flexibility of allowing the
two family lengths to differ, and with almost-sure convergence rather than
merely convergence in probability.  The answering paper's abstract and the
discussion following Theorem~1 explicitly state that Hayes's theorem makes
this a proof of the Peterson--Thom conjecture.  It also notes an independent
Haar-unitary proof by Bordenave and Collins.

\section*{Search scope}

A bounded check through 9 August 2026 used the run registry, exact title and
conjecture wording, tensor-GUE strong convergence, and Peterson--Thom status.
The decisive direct match is arXiv:2205.07695.  Later work, including
arXiv:2308.14109, explicitly calls this the random-matrix solution.  No new
mathematical result is claimed in this packet.

\begin{thebibliography}{9}
\bibitem{source}
B. Hayes,
\emph{A random matrix approach to the Peterson--Thom conjecture},
Indiana Univ. Math. J. \textbf{71} (2022), 1243--1297;
arXiv:2008.12287.

\bibitem{answerpaper}
S. T. Belinschi and M. Capitaine,
\emph{Strong convergence of tensor products of independent G.U.E. matrices},
arXiv:2205.07695, version 2, 2024.
\end{thebibliography}

\end{document}
